\chapter{Non‑Speculative Designs: Single Lock and Serializers}
\index{single global lock}\index{SGL}\index{serializers}\index{queue-based transactional memory}\index{non-speculative transactional memory}\index{opacity}\index{Amdahl's law}

Non-speculative transactional memory is the control case for the rest of the implementation chapters. It gives up parallel execution inside transactions and obtains, in exchange, the simplest correctness story in the book: one transaction runs at a time, observes a sequential state, and publishes the next sequential state.

This chapter studies two forms of that idea. A single-global-lock TM runs the caller inside one process-wide critical section. A thread serializer sends transaction bodies to a single worker and uses the worker's dequeue order as the serialization order. Both are deliberately simple. That simplicity is useful: SGL-like backends in the companion project are reference implementations, regression oracles, and fallbacks for more optimistic systems such as NOrec, TL2, TSC-TM, SwissTM, TinySTM, and the GPU designs.

\section{Single Global Lock TM}

A single-global-lock transactional memory acquires one global mutex at \texttt{tm\_begin} and releases it at \texttt{tm\_commit}. While the mutex is held, ordinary loads and stores are safe because no other transaction can interleave with them. There is no validation phase, no version-clock protocol, no read-set scan, and no rollback of speculative state. The transaction either performs its writes while holding the lock, or performs no writes because its own control flow decides not to.

This design is correct by construction. Every committed transaction has a unique position in the order of lock acquisitions. Every read observes a state produced by the previous committed transaction in that order. Every write becomes visible before the lock is released and before the next transaction begins. Therefore opacity is immediate: even live transactions observe consistent states, because there is never another live writer racing with them.

SGL is also a yardstick. A speculative STM that performs well only on conflict-free benchmarks but collapses into incorrect behavior under contention has failed the basic contract. Under extreme contention, more elaborate systems should degrade toward a safe serial execution, not toward torn reads, failed validation, or memory-management bugs.

The companion codebase exposed one subtle version of this rule. During single-threaded initialization it is tempting to replace \texttt{tm\_malloc}, \texttt{tm\_free}, \texttt{tm\_read}, and \texttt{tm\_write} with direct stubs. That can appear harmless because there is no concurrency yet. It is not harmless. Initialization often creates objects that later participate in transactional validation, allocation tracking, persistence metadata, or ownership protocols. If initialization bypasses the real backend hooks, later tests may execute with silently disabled validation or incomplete metadata. The fixed rule is simple: even single-threaded setup routes memory and access operations through the selected backend.

The PlusCal/SGL model from Part~III is the reference execution: one process owns the lock, executes the transaction body, updates the bank state, and releases the lock. There are no read-write conflicts to discover because the model forbids interleaving at the transaction granularity.

A thread serializer, studied next, is SGL with the mutex replaced by a work queue. It eliminates lock handoff among callers, but it does not create parallel transaction execution. It merely moves the single execution slot from the caller threads to a worker thread.

\begin{figure}[htbp]
\centering
\begin{tikzpicture}[
    node distance=2.2cm,
    state/.style={draw, rounded corners, minimum width=2.8cm, minimum height=0.9cm, align=center},
    edge/.style={-{Latex}, thick}
]
\node[state] (idle) {Idle\\lock free};
\node[state, right=of idle] (active) {Active transaction\\lock held};
\node[state, right=of active] (commit) {Commit\\publish writes};
\node[state, below=of active] (queue) {Serializer queue\\pending tasks};

\draw[edge] (idle) -- node[above]{\texttt{tm\_begin}} (active);
\draw[edge] (active) -- node[above]{\texttt{tm\_commit}} (commit);
\draw[edge] (commit) to[bend left=25] node[below]{unlock} (idle);
\draw[edge] (queue) -- node[right]{dequeue one} (active);
\draw[edge] (idle) -- node[left]{submit} (queue);
\draw[edge] (active) to[loop above] node{reads/writes} (active);
\end{tikzpicture}
\caption{Non-speculative transactional execution. SGL admits one caller by a lock; a serializer admits one submitted task by a queue. Both execute at most one transaction body at a time.}
\label{fig:nonspec-state-machine}
\end{figure}

\begin{table}[htbp]
\centering
\begin{tabular}{M{0.22\linewidth}M{0.23\linewidth}M{0.23\linewidth}M{0.22\linewidth}}
\hline
\thb{0.22\textwidth}{Design} & \thb{0.23\textwidth}{Serialization point} & \thb{0.23\textwidth}{Main cost} & \thb{0.22\textwidth}{Correctness argument} \\
\hline
SGL & Global mutex at \texttt{tm\_begin} & Lock contention and lost parallelism & Mutual exclusion gives a single total order \\
Thread serializer & Worker dequeues one task & Submission, queueing, loss of locality & Worker order is the serialization order \\
DistributedSGL & Distributed lock or sequencer & Network latency and failure handling & One distributed owner at a time \\
PersistentSGL & Persistent critical section protocol & Flushes, fences, recovery metadata & Durable order follows lock-protected commit order \\
\hline
\end{tabular}
\caption{A taxonomy of non-speculative serializers. The mechanism changes, but the proof shape is the same: reduce execution to a sequential history.}
\label{tab:nonspec-taxonomy}
\end{table}

The thread-serializer row is not just a thought experiment: the companion
repository ships one as the \texttt{expli::QueueExecutor}
(\texttt{explicit\_api/cpp/include/\allowbreak tx\_executor.hpp}, backed by
\texttt{queue\_runtime.cpp}), and the bank benchmark runs in this mode with
\texttt{-{}-queue}. Transaction bodies are submitted as closures to a pool of
worker threads; each worker initializes its own TM thread state and drains
its queue serially. Two details repay attention when reading the code: the
worker threads record their stack bounds so the runtime can prove that a
referenced address is stack-local and bypass instrumentation
(\texttt{stm::tm\_record\_\allowbreak stack\_\allowbreak bounds}), and statically allocated
TM-annotated globals are registered explicitly
(\texttt{tm\_register\_\allowbreak global})
because they do not live in the TM region --- both are invisible at the
algorithmic level but load-bearing in practice.

Amdahl's law gives the first-order bound. If fraction f of execution is serialized and N workers are available, the idealized speedup is

$$
S_{\max} = \frac{1}{f + \frac{1-f}{N}}
$$

where f is the non-parallel fraction and N is the number of cores. For SGL, transaction execution itself belongs to the serialized fraction. The exact value depends on the workload, but the shape is invariant: once all useful work is inside the global lock, more cores mostly add waiting.

For example, if nearly all useful application work is performed inside transactions, f approaches one and the speedup approaches one regardless of the number of cores. This is not an implementation accident. It is the price paid for the simple proof.

\section{Thread‑Serializer / Queue‑Based TM}

A thread serializer uses one worker thread to execute submitted transactions sequentially from a queue. The calling threads package transaction bodies as tasks, submit them, and wait for results. The worker repeatedly dequeues one task, runs it to completion, records the result, and moves to the next task.

The important shift is that mutual exclusion is no longer obtained by making many callers contend for a single lock. Instead, callers contend for queue insertion, and the worker owns the mutable state by convention. This can reduce lock handoff costs for tiny transactions, especially when the queue is efficient and the worker remains hot in cache. It can also hurt latency, because every operation now pays submission, wakeup, scheduling, and reply costs.

A serializer is not merely an optimization of SGL. It changes the calling convention. The submitting thread no longer runs the transaction body; the worker does. Therefore transaction-local state cannot be stored only in thread-local variables of the submitter. The descriptor must travel with the task or be rebuilt by the worker.

This matters for real TM backends. A transaction descriptor often contains read sets, write sets, allocation logs, statistics counters, nesting state, epoch information, or persistence metadata. In SGL the descriptor may naturally live on the caller's stack. In a serializer it must be part of the submitted closure, stored in the queue node, or allocated in worker-owned storage. Treating the two designs as interchangeable is a common source of bugs.

Queue serializers also clarify the meaning of a serialization point. For SGL, the serialization point may be taken as the lock acquisition or the commit while the lock is held, depending on the proof style. For a serializer, the serialization order is the dequeue order. The transaction may have been submitted much earlier, but it does not observe memory until the worker runs it.

\section{Worked Example: Bank Transfer Under SGL}
\label{sec:worked-sgl-bank}

The running bank example has two invariants. First, no account may fall below its floor. Second, the total amount of money is conserved by transfers. SGL enforces both by making each transfer execute as one sequential critical section.

The example below keeps the transactional interface even though the implementation is only a mutex. This is intentional. The same application code can later be moved to TL2, NOrec, TinySTM, or another backend without changing every access site. The abstraction boundary is part of the test: reads, writes, allocation, and commit must pass through the backend selected for the run.

\begin{lstlisting}[language=C++, caption={A compilable SGL-style transactional bank transfer.}]
#include <cassert>
#include <cstdint>
#include <iostream>
#include <mutex>
#include <numeric>
#include <stdexcept>
#include <thread>
#include <vector>

struct Account {
    std::int64_t balance;
    std::int64_t floor;
};

class SGLTM {
public:
    class Tx {
    public:
        explicit Tx(SGLTM& tm) : tm_(tm), active_(true) {
            tm_.lock_.lock();
        }

        Tx(const Tx&) = delete;
        Tx& operator=(const Tx&) = delete;

        ~Tx() {
            if (active_) {
                tm_.lock_.unlock();
            }
        }

        template <class T>
        T read(const T& x) const {
            return x;
        }

        template <class T>
        void write(T& x, const T& v) {
            x = v;
        }

        void commit() {
            if (!active_) {
                throw std::logic_error("double commit");
            }
            active_ = false;
            tm_.lock_.unlock();
        }

    private:
        SGLTM& tm_;
        bool active_;
    };

    Tx begin() {
        return Tx(*this);
    }

private:
    std::mutex lock_;
};

class Bank {
public:
    Bank(std::size_t n, std::int64_t initial, std::int64_t floor)
        : accounts_(n, Account{initial, floor}) {}

    bool transfer(std::size_t from, std::size_t to, std::int64_t amount) {
        if (from == to || from >= accounts_.size() || to >= accounts_.size()) {
            return false;
        }

        auto tx = tm_.begin();

        const std::int64_t from_balance = tx.read(accounts_[from].balance);
        const std::int64_t from_floor = tx.read(accounts_[from].floor);
        const std::int64_t to_balance = tx.read(accounts_[to].balance);

        if (from_balance - amount < from_floor) {
            tx.commit();
            return false;
        }

        tx.write(accounts_[from].balance, from_balance - amount);
        tx.write(accounts_[to].balance, to_balance + amount);

        tx.commit();
        return true;
    }

    std::int64_t total() const {
        return std::accumulate(
            accounts_.begin(), accounts_.end(), std::int64_t{0},
            [](std::int64_t s, const Account& a) { return s + a.balance; });
    }

private:
    SGLTM tm_;
    std::vector<Account> accounts_;
};

int main() {
    constexpr std::size_t accounts = 8;
    constexpr std::int64_t initial = 100;
    constexpr std::int64_t floor = 0;

    Bank bank(accounts, initial, floor);
    const std::int64_t before = bank.total();

    std::vector<std::thread> threads;
    for (std::size_t t = 0; t < 4; ++t) {
        threads.emplace_back([&bank, t] {
            for (int i = 0; i < 10000; ++i) {
                const std::size_t from = (t + i) % 8;
                const std::size_t to = (from + 1) % 8;
                bank.transfer(from, to, 1);
            }
        });
    }

    for (auto& th : threads) {
        th.join();
    }

    const std::int64_t after = bank.total();
    assert(before == after);
    std::cout << "money conserved: " << after << "\n";
}
\end{lstlisting}

The linearization point is the interval protected by the global mutex. Reads and writes still go through \texttt{read} and \texttt{write}; replacing them with stubs during initialization can hide validation bugs in later backends. Failed transfers still commit a read-only transaction; there is no rollback path because no speculative write has escaped. The proof is independent of the number of calling threads: all completed transfers form the mutex acquisition order.

The floor invariant is checked before either balance is changed. Since no other transaction can run between the check and the writes, the check remains valid until commit. The conservation invariant follows because the transaction subtracts and adds the same amount in one critical section.

$$
\sum_i balance_i' = \sum_i balance_i
$$

\section{Worked Example: Queue Serializer With Per-Task Descriptors}
\label{sec:worked-serializer-bank}

A serializer moves execution to a worker thread. The transaction descriptor is therefore attached to the submitted task. In this minimal example the descriptor is just the closure itself plus the returned future; real systems attach allocation logs, statistics, or backend-specific metadata.

The bank object below is intentionally not internally synchronized. All access to it is made through the serializer. This is the same proof technique as SGL, but expressed as ownership transfer: while the worker executes a task, it is the only thread touching the bank state.

\begin{lstlisting}[language=C++, caption={A compilable queue-based serializer for bank transfers.}]
#include <cassert>
#include <condition_variable>
#include <cstdint>
#include <future>
#include <iostream>
#include <mutex>
#include <numeric>
#include <queue>
#include <stdexcept>
#include <thread>
#include <type_traits>
#include <vector>
#include <functional>

struct Account {
    std::int64_t balance;
    std::int64_t floor;
};

class Serializer {
public:
    Serializer() : stopping_(false), worker_([this] { run(); }) {}

    Serializer(const Serializer&) = delete;
    Serializer& operator=(const Serializer&) = delete;

    ~Serializer() {
        {
            std::lock_guard<std::mutex> g(m_);
            stopping_ = true;
        }
        cv_.notify_one();
        worker_.join();
    }

    template <class F>
    auto submit(F&& f) -> std::future<std::invoke_result_t<F&>> {
        using R = std::invoke_result_t<F&>;

        auto task = std::make_shared<std::packaged_task<R()>>(
            std::forward<F>(f));
        std::future<R> result = task->get_future();

        {
            std::lock_guard<std::mutex> g(m_);
            if (stopping_) {
                throw std::runtime_error("submit after stop");
            }
            q_.push([task] { (*task)(); });
        }

        cv_.notify_one();
        return result;
    }

private:
    void run() {
        for (;;) {
            std::function<void()> job;

            {
                std::unique_lock<std::mutex> lk(m_);
                cv_.wait(lk, [this] { return stopping_ || !q_.empty(); });

                if (stopping_ && q_.empty()) {
                    return;
                }

                job = std::move(q_.front());
                q_.pop();
            }

            job();
        }
    }

    std::mutex m_;
    std::condition_variable cv_;
    std::queue<std::function<void()>> q_;
    bool stopping_;
    std::thread worker_;
};

class Bank {
public:
    Bank(std::size_t n, std::int64_t initial, std::int64_t floor)
        : accounts_(n, Account{initial, floor}) {}

    bool transfer_serial(std::size_t from, std::size_t to,
                         std::int64_t amount) {
        if (from == to || from >= accounts_.size() || to >= accounts_.size()) {
            return false;
        }

        Account& a = accounts_[from];
        Account& b = accounts_[to];

        if (a.balance - amount < a.floor) {
            return false;
        }

        a.balance -= amount;
        b.balance += amount;
        return true;
    }

    std::int64_t total_serial() const {
        return std::accumulate(
            accounts_.begin(), accounts_.end(), std::int64_t{0},
            [](std::int64_t s, const Account& a) { return s + a.balance; });
    }

private:
    std::vector<Account> accounts_;
};

int main() {
    Serializer ser;
    Bank bank(8, 100, 0);

    const std::int64_t before =
        ser.submit([&bank] { return bank.total_serial(); }).get();

    std::vector<std::future<bool>> replies;
    for (int i = 0; i < 40000; ++i) {
        replies.push_back(ser.submit([&bank, i] {
            const std::size_t from = static_cast<std::size_t>(i % 8);
            const std::size_t to = (from + 1) % 8;
            return bank.transfer_serial(from, to, 1);
        }));
    }

    for (auto& f : replies) {
        (void)f.get();
    }

    const std::int64_t after =
        ser.submit([&bank] { return bank.total_serial(); }).get();

    assert(before == after);
    std::cout << "money conserved: " << after << "\n";
}
\end{lstlisting}

The queue order is the serialization order. The bank object is unsynchronized internally; safety comes from the single worker. The submitter waits on a \texttt{future}, which makes latency visible even when the worker has no lock contention. Thread-local transaction descriptors would be wrong here: the submitting thread and executing thread are different.

The serializer can outperform a contended lock only when the avoided lock traffic is larger than submission and reply overhead. It cannot overcome the fundamental sequential bottleneck. Under a burst of requests, the worker becomes the limiting server and queueing delay dominates.

\paragraph{The companion's version: \texttt{bank --queue}.}
The repository's bank benchmark runs its transfer workload through exactly
this pattern, with two production hardenings worth noticing. The executor is
\texttt{expli::QueueExecutor} constructed with as many workers as queues
(\texttt{QueueExecutor(nb\_threads, nb\_threads)}), installed in a global hook
so that every \texttt{TM<T>::transaction(...)} call site is routed through it
transparently --- inline mode executes the body immediately, queue mode
enqueues it. Synchronous calls enqueue the closure plus an atomic done-flag
and spin-wait; asynchronous calls return a future instead. The hardenings:
first, worker threads record their stack bounds at startup
(\texttt{stm::tm\_record\_stack\_bounds}) so the address classifier can prove
stack-locality for bypass decisions, and TM-annotated globals are registered
explicitly (\texttt{tm\_register\_global}) because static storage is outside
the mmap'd TM region --- omit either and the instrumentation silently stops
tracking real accesses, the Bug-1 pattern of \S\ref{sec:debugging-story}.
Second, all workers park on a start-gate barrier before the measurement
region begins, so throughput numbers measure transactions rather than thread
startup. Running it:

\begin{lstlisting}[caption={The companion bank under queue serialization versus inline execution.}]
$ bank -a 64 -t 4 -d 2000            # inline: callers run TXs directly
Mode: inline
...
$ bank -a 64 -t 4 -d 2000 --queue    # queue: N workers drain N queues
Mode: queue (4 workers, 4 queues)
...
\end{lstlisting}

The instructive observation is that queue mode does not beat inline mode on
this workload: each caller-thread/worker pair adds a handoff, and the money
invariant holds in both modes. The exercise is not to find a winner but to see
the serialization point move from ``caller holds the backend'' to ``worker
dequeues the task'' with the application code unchanged --- which is precisely
the property that makes the serializer a drop-in backend rather than a
rewrite.

\section{Calvin: Static Transactions and Deterministic Execution}
\index{static transaction}\index{Calvin}\index{deterministic execution}\index{sequencer}\index{snapshot isolation}
\label{sec:calvin-static}

The serializer of \S\ref{sec:worked-serializer-bank} executes one transaction body at a time, but the calling threads still build the transaction body \emph{speculatively}: they read shared memory, guess, and hope the guess was right. Calvin takes a different position. A transaction is declared \emph{before} it executes: the read-set, the write-set, and the computation are all fixed up front. The sequencer then assigns a total order, snapshots the declared addresses, and runs each body against its own snapshot. Writes are buffered and published only after the body returns successfully.

The word ``static'' is the distinguishing property. In NOrec, TL2, or TinySTM the runtime discovers the read-set and write-set as the transaction executes --- those sets may be different on every retry. In Calvin, the programmer specifies the access sets at declaration time. That static knowledge has two consequences. First, the runtime can validate every declared read against the snapshot \emph{before} executing the body, so no work is wasted on a transaction whose snapshot is already stale. Second, execution is deterministic: given the same snapshot, the body always produces the same writes. Determinism simplifies reasoning about correctness, because the commit of one transaction cannot depend on timing details that differ between runs.

The trade-off is inflexibility. If the read-set is not known until the body runs, the programmer must over-approximate (declaring more addresses than needed, which widens the conflict window) or structure the code so that the access set is visible at the call site. This is natural for message-queue workers, financial settlements, and event-driven simulations where each event's inputs are known before the computation starts, but awkward for interactive code where the read-set depends on data read during execution.

The companion codebase's \texttt{calvin} backend is named after, but is not, the deterministic distributed database of Thomson et al.~\cite{thomson12calvin}; it is a two-phase OCC shared-memory backend. The model below is closer to the original Calvin idea: static declarations, sequenced execution, deterministic results.

\subsection{State: declarations, memory, and a sequencer}

Listing~\ref{lst:calvin-state} defines the shared memory (one atomic word per account), the transaction declaration type (read-set, write-set, and a body that receives a snapshot and returns writes or abort), and the sequencer that dequeues declarations in order, snapshots the read-set, and executes each body.

\begin{lstlisting}[language=Rust, caption={Calvin state: word-addressed memory, static transaction declarations, and the sequencer.}, label={lst:calvin-state}]
use std::sync::atomic::{AtomicU64, AtomicUsize, Ordering as O};
use std::sync::Arc;
use std::thread;

struct Mem { data: Vec<AtomicU64> }
impl Mem {
    fn new(n: usize) -> Self {
        Mem { data: (0..n).map(|_| AtomicU64::new(0)).collect() }
    }
    fn alloc(&mut self, v: u64) -> usize {
        let i = self.data.len();
        self.data.push(AtomicU64::new(v));
        i
    }
    fn get(&self, a: usize) -> u64 {
        self.data[a].load(O::Acquire)
    }
    fn set(&self, a: usize, v: u64) {
        self.data[a].store(v, O::Release);
    }
}

type BodyFn =
    Box<dyn Fn(&[u64]) -> Option<Vec<(usize, u64)>> + Send>;

struct TxDecl {
    reads:  Vec<usize>,            // addresses to snapshot
    writes: Vec<usize>,            // addresses the body may modify
    body:   BodyFn,                // None = abort; Some(writes) = commit
}

struct Sequencer { mem: Arc<Mem>, seqno: AtomicUsize }
impl Sequencer {
    fn new(mem: Arc<Mem>) -> Self {
        Sequencer { mem, seqno: AtomicUsize::new(0) }
    }
    fn snapshot(&self, addrs: &[usize]) -> Vec<u64> {
        addrs.iter().map(|&a| self.mem.get(a)).collect()
    }
\end{lstlisting}

The \texttt{TxDecl} type is the static contract. The \texttt{reads} and \texttt{writes} vectors declare which addresses the body will touch; the \texttt{body} closure receives the snapshot values at those addresses and returns \texttt{None} (abort) or \texttt{Some(writes)} (commit). The body never touches shared memory directly --- all reads come from the snapshot, all writes are buffered.

\subsection{Execute and publish}

Listing~\ref{lst:calvin-exec} is the execution loop. The sequencer dequeues one declaration, snapshots the read-set, calls the body, and publishes the writes. The important property is that the body is never retried: if the snapshot is stale, the declaration fails immediately at the \emph{declaration} phase, not the \emph{execution} phase. This is the fundamental difference from speculative STMs: there is no ``abort and retry the body'' because the body is never allowed to see stale data.

\begin{lstlisting}[language=Rust, caption={Calvin execution: snapshot, run body, publish writes. No retry needed.}, label={lst:calvin-exec}]
    fn run(&self, queue: Vec<TxDecl>) {
        for decl in queue {
            let seq = self.seqno.fetch_add(1, O::Relaxed);
            let snap = self.snapshot(&decl.reads);
            match (decl.body)(&snap) {
                None => {
                    eprintln!("[seq={}] abort (stale snap)", seq);
                }
                Some(writes) => {
                    for (a, v) in writes {
                        self.mem.set(a, v);
                    }
                    println!("[seq={}] committed  snap={:?}",
                             seq, snap);
                }
            }
        }
    }
}
\end{lstlisting}

The sequencer is single-threaded, which makes the total order explicit: the dequeue order \emph{is} the serialization order. In a distributed Calvin system the sequencer would be a separate process assigning sequence numbers; here the in-process loop is enough to show the idea. The \texttt{AtomicUsize} counter records how many transactions have been processed, and the \texttt{Relaxed} ordering on \texttt{fetch\_add} is sufficient because the total order is established by the sequential loop itself, not by the counter value.

\subsection{Worked Example: Bank Transfer}

Listing~\ref{lst:calvin-demo} builds the transaction queue: transfer declarations interleaved with read-only checker assertions. The deterministic Fisher-Yates shuffle (fixed seed) ensures reproducible interleaving across runs. Each transfer declares its read-set and write-set explicitly; the checker declares read-only addresses and asserts conservation.

\begin{lstlisting}[language=Rust, caption={Calvin bank demo: static transfer declarations, deterministic shuffle, conservation assertion mid-flight.}, label={lst:calvin-demo}]
const ITERS: u64 = 500;
const START: u64 = 1_000_000;

fn main() {
    let mut mem = Mem::new(0);
    let a = mem.alloc(START);
    let b = mem.alloc(START);
    let mem = Arc::new(mem);

    let mut txs: Vec<TxDecl> = Vec::new();
    for t in 0..8 {
        let (src, dst) = if t % 2 == 0 { (a, b) }
                         else { (b, a) };
        for _ in 0..ITERS {
            txs.push(TxDecl {
                reads:  vec![src, dst],
                writes: vec![src, dst],
                body: Box::new(move |snap| {
                    let (s, d) = (snap[0], snap[1]);
                    if s < 10 { return None; }
                    Some(vec![(src, s - 10), (dst, d + 10)])
                }),
            });
        }
    }
    for _ in 0..500 {
        txs.push(TxDecl {
            reads:  vec![a, b],
            writes: vec![],
            body: Box::new(move |snap| {
                assert_eq!(snap[0] + snap[1], 2 * START);
                Some(vec![])
            }),
        });
    }
    // Deterministic shuffle (fixed seed).
    { let mut s: u64 = 0xDEAD_BEEF_CAFE_1234;
      for i in (1..txs.len()).rev() {
          s = s.wrapping_mul(6364136223846793005).wrapping_add(1);
          let j = (s >> 33) as usize % (i + 1);
          txs.swap(i, j);
      }
    }

    let seq = Sequencer::new(mem.clone());
    thread::spawn(move || seq.run(txs)).join().unwrap();

    let (fa, fb) = (mem.get(a), mem.get(b));
    println!("\nfinal A={} B={} total={}", fa, fb, fa + fb);
    assert_eq!(fa + fb, 2 * START);
}
\end{lstlisting}

The checker runs inside the sequencer's total order, interleaved with transfers at deterministic positions. Because the snapshot is taken before each body runs, the checker sees a consistent world at every point --- money is either pre-transfer or post-transfer, never half-transferred. This is opacity by construction: no concurrent mutation can sneak between the snapshot and the body because the sequencer is single-threaded.

Three properties distinguish Calvin from the speculative designs of Chapter~9:

\begin{itemize}
    \item \textbf{No abort-retry loop}. A speculative STM retries the body when validation fails; Calvin rejects stale declarations before the body runs. The body is never wasted work.
    \item \textbf{Determinism}. Given the same declaration queue, the commit order and the final memory state are identical across runs. Speculative STMs can produce different outcomes depending on thread scheduling.
    \item \textbf{Static access sets}. The read-set and write-set are visible at declaration time, enabling the runtime to batch, reorder, or optimize declarations before executing them.
\end{itemize}

The price is the loss of dynamic read-sets. A body that reads an address discovered during execution cannot declare that address up front, so Calvin is best suited for workloads where inputs are known: message processing, event handlers, map-reduce workers, and batch settlement. For interactive or exploratory code, the speculative designs of Chapter~9 remain the right tool.

\section{Summary}

Non-speculative designs are correct by construction and serve as a baseline. Their lack of concurrency limits scalability, but their proof obligations are small: every transaction observes a state produced by the previous committed transaction.

SGL is the simplest opacity proof in the book. Queue serializers replace lock handoff with task handoff. Calvin replaces speculative execution with static declarations and deterministic sequencing: the read-set is fixed before the body runs, so stale snapshots are rejected without wasting execution work. The trade-off is inflexibility --- the access set must be known at declaration time.

In the companion codebase, SGL-like designs are not toys. They serve as reference backends and regression oracles alongside speculative systems such as NOrec, TL2, TSC-TM, SwissTM, TinySTM, and others. The reported test suites pass \texttt{test\_tx} 114/114 and \texttt{test\_ds} 207/207 for the implemented backends. Measured results on the M1 Pro fuzz/bank workload also show why the baseline matters: TL2 and TSC-TM reach about 2.06M and 2.05M transactions per second in the reported runs, while NOrec-BF improves over NOrec on the read-mostly case. These numbers are not claims about this chapter's designs; they are evidence that the simple serial yardstick remains useful when evaluating faster speculative backends.

\section{Exercises}
\begin{enumerate}
    \item Calculate the maximum speedup achievable with a single‑global‑lock TM on a 16‑core machine when transactions spend 10\% of their time in a critical section.
    \item A serializing TM uses a single worker thread. Explain why this design can still outperform fine‑grained locking if transactions are very short and lock overhead dominates.
    \item Propose a hybrid design that starts with a non‑speculative execution and switches to speculative when contention is low. What information does the runtime need to make the switch?
    \item \emph{[implementation]} Extend Listing~\ref{sec:worked-sgl-bank} with \texttt{tm\_malloc} and \texttt{tm\_free} hooks. The hooks may still use the process allocator, but calls must pass through the backend object. Explain which tests would fail to detect bugs if initialization bypassed these hooks.
    \item \emph{[verification]} Write a PlusCal model of the bank transfer under SGL with two accounts, one lock, and two clients. State the money conservation invariant and the account-floor invariant. Model-check both with all transfer amounts in the set \texttt{\{0,1,2\}}.
    \item \emph{[theory]} Compare SGL and a queue serializer with three costs: lock acquisition, queue submission, and transaction body execution. Derive the condition under which the serializer has lower latency for a single transaction, and then explain why that condition may not imply higher throughput under a burst of requests.
    \item \emph{[implementation]} Modify the queue serializer so that each task carries an explicit transaction descriptor containing a task id, an allocation log, and a commit flag. Ensure that the descriptor is created by the submitter but finalized by the worker.
    \item \emph{[verification]} State the refinement mapping from a queue serializer to an abstract sequential TM. Identify the concrete variable that determines the abstract transaction order.
    \item \emph{[theory]} For a workload in which every transaction runs under SGL, explain why increasing the number of caller threads can reduce throughput even though it does not change the sequential work performed.
    \item \emph{[implementation]} Build the Calvin bank demo of \S\ref{sec:calvin-static} from the source in Appendix~\ref{chap:rust-mistakes}. Replace the deterministic shuffle with a random one and run it five times. Verify that the final total is always \texttt{2000000} but the committed snapshot values differ. Explain why determinism of the commit \emph{order} does not affect correctness here.
    \item \emph{[theory]} Calvin rejects stale declarations before executing the body. A speculative STM (Chapter~9) executes the body first, then validates. Derive the condition under which Calvin does strictly less work than a speculative STM for a given workload, and the condition under which the opposite holds. Hint: compare the cost of body execution with the cost of snapshot validation.
\end{enumerate}

\textbf{How to check your work.} The \emph{[implementation]} exercises are
verified with the companion: route \texttt{tm\_malloc}/\texttt{tm\_free}
through the backend of \S\ref{sec:worked-sgl-bank}, and finalize the
serializer task descriptor per \S\ref{sec:worked-serializer-bank}.  The
\emph{[verification]} exercises are checked with TLC on the SGL bank model
(\texttt{MoneyConserved}, account floor; see
Appendix~\ref{chap:pluscal-quick-ref}).  The \emph{[theory]} exercises are
graded against \S\ref{sec:worked-sgl-bank} and \S\ref{sec:semantic-obligations}.

\index{global mutex}
\index{transaction descriptor}
\index{queue submission}
\index{worker thread}
\index{serialization order}
\index{lock acquisition order}
\index{linearization point}
\index{bank transfer invariant}
\index{money conservation}
\index{account floor invariant}
\index{backend hooks}
\index{initialization stubs}
\index{thread-local state}
\index{task handoff}
\index{regression oracle}
\index{DistributedSGL}
\index{PersistentSGL}

% -------------------------------------------------------------------
